\section*{Цели работы}
Практическое ознакомление с принципами реализации и влиянием различных типов отрицательной обратной связи (ООС) на характеристики усилительных каскадов.

\section*{Задачи}
\begin{itemize}
    \item Исследование амплитудно-частотных характеристик (АЧХ) усилителя при четырех типах ООС.
    \item Определение граничных частот и полосы пропускания для каждой конфигурации.
    \item Экспериментальное определение входного и выходного сопротивлений усилителя при наличии ООС.
\end{itemize}

\section*{Краткие теоретические сведения}
Обратная связь (ОС) — это передача части энергии выходного сигнала на вход системы. В данной работе исследуется отрицательная ОС, при которой сигнал ОС подается в противофазе с входным. Основные параметры:
\begin{itemize}
    \item Коэффициент обратной связи: $\gamma = U_{o.c} / U_{вых}$.
    \item Коэффициент усиления с ОС:
    \begin{equation*}
        K_{o.c} = \frac{K}{1 + \gamma K}
    \end{equation*}
    где $K$ — коэффициент усиления без ОС.
\end{itemize}
Классификация ОС осуществляется по способу снятия сигнала с выхода (по напряжению или по току) и по способу введения на вход (последовательная или параллельная). ООС стабилизирует коэффициент усиления, расширяет полосу пропускания, а также изменяет входное и выходное сопротивления в $(1 + \gamma K)$ раз.

\section*{Инструкция по снятию данных}

\subsection*{1. Подготовка стенда}
\begin{itemize}
    \item \textit{Что сделать}: Подать напряжение питания на макет.
    \item \textit{Как сделать}: Соединить клеммы источника постоянного тока с гнездами «$+15$ В» и «$-15$ В» на лицевой панели макета.
    \item \textit{Нюанс}: Перед включением источника убедитесь в правильности полярности подключения, чтобы избежать выхода операционных усилителей из строя.
\end{itemize}

\subsection*{2. Исследование АЧХ}
\begin{itemize}
    \item \textit{Что сделать}: Снять зависимость выходного напряжения от частоты для четырех конфигураций ООС.
    \item \textit{Как сделать}: 
    \begin{enumerate}
        \item Соберите схему с параллельной ООС по напряжению с помощью перемычек.
        \item Подайте на вход сигнал от генератора гармонических колебаний. Установите амплитуду $U_{вх}$ так, чтобы на средних частотах отсутствовали видимые искажения на осциллографе.
        \item Изменяйте частоту генератора в диапазоне от 20 Гц до 3 МГц согласно точкам в таблице.
        \item Фиксируйте значения $U_{вых}$ по вольтметру.
        \item Повторите измерения для схем: последовательная по напряжению, параллельная по току и последовательная по току.
    \end{enumerate}
    \item \textit{Предупреждение}: Поддерживайте амплитуду входного сигнала $U_{вх}$ постоянной во всем диапазоне частот.
\end{itemize}

\subsection*{3. Определение входного сопротивления}
\begin{itemize}
    \item \textit{Что сделать}: Измерить $R_{вх}$ для параллельной и последовательной ООС.
    \item \textit{Как сделать}: Включите в разрыв входной цепи магазин сопротивлений. Подберите такое сопротивление $R_{маг}$, при котором выходное напряжение уменьшится вдвое по сравнению с исходным. В этом случае $R_{вх} = R_{маг}$. Данные занесите в таблицу.
    \item \textit{Нюанс}: Измерения проводить на частоте 1 кГц (область средних частот).
\end{itemize}

\subsection*{4. Определение выходного сопротивления}
\begin{itemize}
    \item \textit{Что сделать}: Измерить $R_{вых}$ для ООС по току и по напряжению.
    \item \textit{Как сделать}: Измерьте $U_{вых}$ на холостом ходу ($R_H = \infty$), затем подключите известную нагрузку $R_H$ и измерьте $U_{вых.н}$. Вычислите $R_{вых}$ по формуле падения напряжения. Повторите для разных типов ОС согласно таблице.
    \item \textit{Нюанс}: При измерении ООС по току следите, чтобы цепь нагрузки не разрывалась полностью, если это предусмотрено конструкцией цепи ОС.
\end{itemize}

% \begin{figure}[H]
%     \centering
%     \includegraphics[width=0.6\textwidth]{figs/fig5_1.pdf}
%     \caption{Схема с последовательной ООС по напряжению}
%     \label{fig:fig5_1}
% \end{figure}

% \begin{figure}[H]
%     \centering
%     \includegraphics[width=0.6\textwidth]{figs/fig5_2.pdf}
%     \caption{Схема с параллельной ООС по напряжению}
%     \label{fig:fig5_2}
% </figure>

% \begin{figure}[H]
%     \centering
%     \includegraphics[width=0.6\textwidth]{figs/fig5_3.pdf}
%     \caption{Схема с последовательной ООС по току}
%     \label{fig:fig5_3}
% \end{figure}

% \begin{figure}[H]
%     \centering
%     \includegraphics[width=0.6\textwidth]{figs/fig5_4.pdf}
%     \caption{Схема с параллельной ООС по току}
%     \label{fig:fig5_4}
% \end{figure}

% \begin{figure}[H]
%     \centering
%     \includegraphics[width=0.8\textwidth]{figs/fig5_5.pdf}
%     \caption{Лицевая панель лабораторного макета}
%     \label{fig:fig5_5}
% \end{figure}

% Соответствие изображений:
% fig:fig5_1 -> Рис. 5.1 (стр. 29)
% fig:fig5_2 -> Рис. 5.2 (стр. 29)
% fig:fig5_3 -> Рис. 5.3 (стр. 30)
% fig:fig5_4 -> Рис. 5.4 (стр. 30)
% fig:fig5_5 -> Рис. 5.5 (стр. 31)